\documentclass[a4paper,11pt,lualatex,ja=standard]{bxjsarticle}

% 数式
\usepackage{amsmath,amsfonts}
\usepackage{bm}
% 画像
\usepackage{graphicx}
% url
\usepackage{hyperref}

% ===== cover command =====
\newcommand{\cover}[3]{
  \thispagestyle{empty}
  \vspace*{4cm}
  \begin{center}
    {\LARGE #1}\\[2cm]
    {\Large #2}\\[3cm]
    {\large 氏名：#3}\\[1cm]
    % {\large 学籍番号：s2510063}
  \end{center}
  \newpage
}
% =========================

\begin{document}

% ===== cover =====
\cover{bytemancy 1}{CyLab}{yukimi}
% =================

\title{レポートタイトル}
\author{氏名}
\date{\today}

\section{問題}
\begin{verbatim}
  while(True):
  try:
    print('⊹──────[ BYTEMANCY-1 ]──────⊹')
    print("☍⟐☉⟊☽☈⟁⧋⟡☍⟐☉⟊☽☈⟁⧋⟡☍⟐☉⟊☽☈⟁⧋⟡☍⟐")
    print()
    print('Send me ASCII DECIMAL 101 1751 times, side-by-side, no space.')
    print()
    print("☍⟐☉⟊☽☈⟁⧋⟡☍⟐☉⟊☽☈⟁⧋⟡☍⟐☉⟊☽☈⟁⧋⟡☍⟐")
    print('⊹─────────────⟡─────────────⊹')
    user_input = input('==> ')
    if user_input == "\x65"*1751:
      print(open("./flag.txt", "r").read())
      break
    else:
      print("That wasn't it. I got: " + str(user_input))
      print()
      print()
      print()
  except Exception as e:
    print(e)
    break
\end{verbatim}

\section{回答}
\begin{verbatim}
yukimi-academy@webshell:~$  nc foggy-cliff.picoctf.net 60722
⊹──────[ BYTEMANCY-1 ]──────⊹
☍⟐☉⟊☽☈⟁⧋⟡☍⟐☉⟊☽☈⟁⧋⟡☍⟐☉⟊☽☈⟁⧋⟡☍⟐

Send me ASCII DECIMAL 101 1751 times, side-by-side, no space.

☍⟐☉⟊☽☈⟁⧋⟡☍⟐☉⟊☽☈⟁⧋⟡☍⟐☉⟊☽☈⟁⧋⟡☍⟐
⊹─────────────⟡─────────────⊹
==> eeeeeeeeeeeeeeeeeeeeeeeeeeeeeeeeeeeeeeeeeeeeeeeeeeeeeeeeeeeeeeeeeeeeeeeeeeeeeeeeeeeeeeeeeeeeeeeeeeeeeeeeeeeeeeeeeeeeeeeeeeeeeeeeeeeeeeeeeeeeeeeeeeeeeeeeeeeeeeeeeeeeeeeeeeeeeeeeeeeeeeeeeeeeeeeeeeeeeeeeeeeeeeeeeeeeeeeeeeeeeeeeeeeeeeeeeeeeeeeeeeeeeeeeeeeeeeeeeeeeeeeeeeeeeeeeeeeeeeeeeeeeeeeeeeeeeeeeeeeeeeeeeeeeeeeeeeeeeeeeeeeeeeeeeeeeeeeeeeeeeeeeeeeeeeeeeeeeeeeeeeeeeeeeeeeeeeeeeeeeeeeeeeeeeeeeeeeeeeeeeeeeeeeeeeeeeeeeeeeeeeeeeeeeeeeeeeeeeeeeeeeeeeeeeeeeeeeeeeeeeeeeeeeeeeeeeeeeeeeeeeeeeeeeeeeeeeeeeeeeeeeeeeeeeeeeeeeeeeeeeeeeeeeeeeeeeeeeeeeeeeeeeeeeeeeeeeeeeeeeeeeeeeeeeeeeeeeeeeeeeeeeeeeeeeeeeeeeeeeeeeeeeeeeeeeeeeeeeeeeeeeeeeeeeeeeeeeeeeeeeeeeeeeeeeeeeeeeeeeeeeeeeeeeeeeeeeeeeeeeeeeeeeeeeeeeeeeeeeeeeeeeeeeeeeeeeeeeeeeeeeeeeeeeeeeeeeeeeeeeeeeeeeeeeeeeeeeeeeeeeeeeeeeeeeeeeeeeeeeeeeeeeeeeeeeeeeeeeeeeeeeeeeeeeeeeeeeeeeeeeeeeeeeeeeeeeeeeeeeeeeeeeeeeeeeeeeeeeeeeeeeeeeeeeeeeeeeeeeeeeeeeeeeeeeeeeeeeeeeeeeeeeeeeeeeeeeeeeeeeeeeeeeeeeeeeeeeeeeeeeeeeeeeeeeeeeeeeeeeeeeeeeeeeeeeeeeeeeeeeeeeeeeeeeeeeeeeeeeeeeeeeeeeeeeeeeeeeeeeeeeeeeeeeeeeeeeeeeeeeeeeeeeeeeeeeeeeeeeeeeeeeeeeeeeeeeeeeeeeeeeeeeeeeeeeeeeeeeeeeeeeeeeeeeeeeeeeeeeeeeeeeeeeeeeeeeeeeeeeeeeeeeeeeeeeeeeeeeeeeeeeeeeeeeeeeeeeeeeeeeeeeeeeeeeeeeeeeeeeeeeeeeeeeeeeeeeeeeeeeeeeeeeeeeeeeeeeeeeeeeeeeeeeeeeeeeeeeeeeeeeeeeeeeeeeeeeeeeeeeeeeeeeeeeeeeeeeeeeeeeeeeeeeeeeeeeeeeeeeeeeeeeeeeeeeeeeeeeeeeeeeeeeeeeeeeeeeeeeeeeeeeeeeeeeeeeeeeeeeeeeeeeeeeeeeeeeeeeeeeeeeeeeeeeeeeeeeeeeeeeeeeeeeeeeeeeeeeeeeeeeeeeeeeeeeeeeeeeeeeeeeeeeeeeeeeeeeeeeeeeeeeeeeeeeeeeeeeeeeeeeeeeeeeeeeeeeeeeeeeeeeeeeeeeeeeeeeeeeeeeeeeeeeeeeeeeeeeeeeeeeeeeeeeeeeeeeeeeeeeeeeeeeeeeeeeeeeeeeeeeeeeeeeeeeeeeeeeeeeeeeeeeeeeeeeeeeeeeeeeeeeeeeeeeeeeeeeeeeeeeeeeeeeeeeeeeeeeeeeeeeeeeeeeeeeeeeeeeeeeeeeeeeeeeeeeeeeeeeeeeeeeeeeeeeeeeeeeeeeeeeeeeeeeeee
picoCTF{h0w_m4ny_e's???_446bf6f1}



\end{verbatim}

\section{メモ}
\begin{verbatim}
  ASCII DECIMAL 101q
  ASCIIコードの10進数表記で101
  
  "\x65"*1751 := e*1751
  
  python3 -c 'print("e"*1751)'と"e"*1751の違い
  
  "e"*1751
→ Pythonの式そのもの
→ 入力欄に打つと、ただの文字列として扱われる
→ だから失敗する

python3 -c 'print("e"*1751)'
→ Pythonに式を実行させるコマンド
→ 結果として eが1751個 出力される
→ その出力を入力すれば成功する

>>> print('e'*1751)でeを1751個並べた

\end{verbatim}

\end{document}